%% introduction.tex

\section{Introduction}
\label{sec:introduction}

\paragraph{Context}
Modern LLM agent systems increasingly adopt operating-system design principles: agents are scheduled as processes, tenants are isolated, tool invocations are gated by permission checkers, and---crucially---\emph{automated self-healing} is employed to handle the frequent failures inherent to LLM execution~\cite{autogen,swe-agent,omo}. When an agent fails (empty response, tool error, verification failure, or semantic crash), a recovery orchestrator classifies the fault and dispatches a layered chain of recovery strategies: re-injection of error context into the next prompt, model fallback, topology mutation, and human escalation. This design is motivated by the empirical observation that LLM agents fail often, and that retrying with structured diagnostic context frequently succeeds where the original attempt did not. As a result, self-healing has become a load-bearing component of production agent frameworks.

\paragraph{Gap}
Despite the richness of this recovery machinery, its security properties have not been examined. The recovery pipeline inherits an implicit assumption from classical OS exception handling: \emph{that failure is a benign accident and that retry is safe.} This assumption holds in classical kernels because exception payloads (e.g., page-fault addresses) are generated by the trusted computing base. It \emph{does not hold} in LLM agent systems, where the failure payload is frequently \emph{external content that the agent was processing when it crashed}. An agent that fetches an adversarial web page, throws an exception whose trace embeds the page text, and hands that trace to the recovery pipeline is effectively hand-delivering attacker-controlled content to a path that was never designed to receive untrusted input. The security literature has extensively studied prompt injection on the \emph{forward} execution path~\cite{greshake2023indirect,liu2024formalizing}, but the \emph{reverse} recovery path---from failure, through diagnostic logging, back into the next-round LLM context---has remained an unexamined trust boundary.

\paragraph{Problem}
We introduce \textbf{Recovery-Channel Injection (RCI)}, a class of attacks that exploit this gap. The key distinction from conventional prompt injection is the \emph{trust-escalation trampoline}: content that was correctly quarantined as untrusted on the forward path is re-laundered as trusted on the recovery path. Conventional prompt injection flows along the forward path---attacker-controlled input enters the LLM context directly, and defenses (input sanitization, tool-permission checks, instruction-hierarchy enforcement) are applied at the entry point~\cite{instruction-hierarchy,spotlighting}. Recovery-channel injection flows along the reverse path: adversarial content enters as legitimate external data, \emph{causes a failure}, and then re-enters the LLM context via the recovery pipeline---a path that the entry-point defenses do not cover. Critically, the recovery pipeline \emph{actively upgrades} the trust level of the content it processes, wrapping error traces in system-level framing and executing recovery actions without re-validation. No amount of forward-path hardening closes this, because the bypass occurs \emph{after} the forward path has legitimately accepted the content.

We instantiate this attack surface in \system{}, a multi-agent operating system prototype developed by the authors with a multi-tier recovery architecture (see \S\ref{sec:threat-model}), and study two structural attack instances that concretely realize the trampoline:

\begin{itemize}[leftmargin=*]
    \item \textbf{(V1) Trust escalation via high-trust template re-injection (\code{ReflectionInjectionRecovery}).} This strategy injects the previous failure's error trace into the next-round prompt under a \code{[SYSTEM CRITICAL - PREVIOUS ATTEMPT FAILED]} template. The trust origin of the trace is resolved via \code{TrustOrigin.fromMetadata()}, which \emph{defaults to \code{SYSTEM\_GENERATED}} when the upstream capture site omits a provenance tag. Consequently, externally-originated traces---e.g., from \code{web\_fetch} on an adversarial page---are silently treated as system-trusted and framed as privileged directives. Compounding this, the \code{RecoveryPromptSanitizer} neutralizes only structural control tokens (\code{<tool\_call>}) and is transparent to natural-language injection directives. In our real-LLM attack harness, a canary-tool directive embedded in an external page triggered in \textbf{57\%} of baseline trials (with the median being 100\%).

    \item \textbf{(V2) Privilege escalation via zero-permission topology mutation (\code{TopologyMutationStrategy}).} This strategy reads a semantic core dump, queries an LLM for a capability-mismatch diagnosis, and passes the LLM-suggested replacement role directly to \code{WorkflowEngine.resumeNode()}---the primitive that swaps a crashed node's role---with \emph{no permission re-check}. Because the core dump captures the agent's last input (which may be external content), an attacker can embed text that induces the diagnostic LLM to emit a high-privilege role name (e.g., \code{System\_Admin}). In baseline configuration, this attacker-suggested role reached the role-replacement primitive in \textbf{100\%} of trials. We emphasize that V2 is an instance of CWE-862 (Missing Authorization)~\cite{cwe862}---a decades-old vulnerability class---not a novel primitive. Its contribution is \emph{not} the discovery of CWE-862, but the measurement of \emph{where, how, and why} this classic vulnerability systematically appears in self-healing agent architectures: recovery actions bypass the permission checks enforced on the forward path because the recovery pipeline is assumed to be benign, and the confused-deputy pattern~\cite{hardy1988confused} arises naturally when a privileged recovery deputy consumes attacker-influenced content. Framing V2 this way makes the contribution empirical (``CWE-862 manifests systematically in self-healing recovery paths, and here is why'') rather than a claim of novelty.
\end{itemize}

\paragraph{Contributions}
We employ a two-stage methodology common in security research: a \emph{white-box testbed} for mechanism attribution and a \emph{black-box production framework} for external validity. Because production-grade agent frameworks do not expose sufficient internal hooks for fine-grained mechanism attribution (e.g., isolating whether content-level provenance or structural privilege enforcement drives a given outcome), we construct \system{} as a white-box testbed; we then validate on LangGraph, a production framework not developed by the authors, to establish external validity. This paper makes the following contributions:

\begin{itemize}[leftmargin=*]
    \item \textbf{A new attack surface and formalization.} We define and characterize \emph{Recovery-Channel Injection (RCI)}, distinguishing it from conventional prompt injection by the trust-escalation trampoline it exploits. We identify the systemic assumption (``failure = benign, retry = safe'') that creates the surface, and formalize the trust gap between forward-path and reverse-path content handling.

    \item \textbf{White-box testbed mechanism attribution (\system{}).} We construct \system{}, a multi-agent operating system prototype with a multi-tier recovery architecture, as a white-box testbed to isolate and quantify two attack mechanisms: (V1) trust escalation via \code{ReflectionInjectionRecovery} and (V2) privilege escalation via \code{TopologyMutationStrategy}. We emphasize that V2 is not a novel primitive---it is CWE-862 (Missing Authorization)~\cite{cwe862}, a decades-old vulnerability class---but its \emph{systematic appearance} in self-healing recovery paths (where the recovery pipeline is assumed benign and bypasses forward-path permission checks) is an empirical finding. The testbed's role is mechanism attribution: isolating content-level (V1) vs.\ structural-level (V2) contributions. In baseline configuration, V2 achieves 100\% ASR (Wilson 95\% CI $[0.96, 1.00]$) and V1 57\% ASR ($[0.47, 0.66]$) against real LLMs. The proposed defense (\trustorigin{} Tagging + \reauthgate{}) reduces V1 to 0.07 ($[0.03, 0.14]$; $p < 10^{-6}$) and V2 to 0.00 (deterministic), with 0.0 FAR and sub-millisecond overhead.

    \item \textbf{External validity on LangGraph (production framework).} The testbed establishes the \emph{mechanism}; LangGraph establishes that the vulnerability class and defense principle are not artifacts of the authors' prototype. A forward-path experiment on \code{create\_react\_agent} using entirely library-default components (no author-constructed code) achieves 0.88 V1 pooled ASR ($[0.76, 0.94]$), confirming that production frameworks lack provenance tagging---though this is the classic indirect injection vector~\cite{greshake2023indirect}, not RCI per se. A recovery-path experiment on \code{StateGraph} using an author-constructed recovery node (following common Reflexion-style patterns) achieves 0.76 pooled ASR ($[0.57, 0.89]$), reduced to 0.04 ($[0.01, 0.20]$) by the untrusted-framing defense, demonstrating the RCI mechanism is realizable on a third-party API; V2 achieves 0.54 ($[0.40, 0.67]$), reduced to 0.00 by the \reauthgate{} analog.

    \item \textbf{Critical experiment: RCI is an independent attack surface (\S\ref{sec:forward-defense-bypass}).} To establish that RCI is not merely a variant of indirect prompt injection, we run the V1 attack with three forward-path defenses fully enabled (input sanitizer + LLM injection detector + tool-permission check). The attack still achieves 0.68 ASR ($[0.54, 0.79]$), statistically indistinguishable from the undefended 0.74 ($p=0.660$). Forward-path defenses blocked 0/50 trials because they inspect user input (benign), while the adversarial payload enters via the tool's failure---which is never inspected. This confirms that RCI bypasses forward-path defenses by construction: the recovery pipeline re-stamps untrusted content as system-trusted \emph{after} those defenses have already accepted it.

    \item \textbf{Wild-case study on production agent codebases with confirmed end-to-end exploit.} To address the concern that the LangGraph recovery node is author-constructed (LangGraph ships no built-in recovery primitive), we conducted a source-level wild-case study of four externally-maintained, production-deployed agent codebases---Reflexion~\cite{reflexion,reflexion-code}, MetaGPT~\cite{meta-agent}, Aider~\cite{aider-code}, OpenHands~\cite{openhands-code}---with a combined $\sim$105k GitHub stars. Three of four ship retry/reflection logic with the structural signature of the V1 trust-escalation trampoline (no provenance tag $\wedge$ high-trust frame). For Reflexion, we additionally executed a \emph{confirmed end-to-end exploit} against its actual \code{generic\_generate\_func\_impl} retry path (commit \texttt{218cf0e}): an adversarial test case whose output contains an injection directive is captured as \code{feedback}, wrapped in a \code{[unit test results]} envelope under a system-framed prompt, and delivered to the LLM---which obeys and emits \code{import canary\_probe; canary\_probe()} in its ``improved implementation''. The exploit achieves 0.52 pooled ASR ($[0.39, 0.65]$, $p = 3.5 \times 10^{-10}$ vs.\ 0\% FPR benign control), with two of five variants at 100\% ASR. This uses Reflexion's verbatim prompt construction and \code{parse\_code\_block}---no author-constructed recovery node.

    \item \textbf{Stronger-model evaluation revealing a content-vs-structural defense asymmetry.} We evaluate on gpt-oss-120b (OpenAI's 120B open-weights model, 200 additional API calls; not a closed-weight frontier model) and find that content-level V1 defense \emph{degrades} on stronger instruction-following models (defended ASR: 0.07 on DeepSeek/Qwen $\rightarrow$ 0.50 on gpt-oss-120b), while structural V2 defense remains at 0.00 ASR across all models. This suggests that as models become more instruction-following, structural privilege enforcement transitions from complementary to \emph{necessary}. A supplementary evaluation on a closed-weight model (accessed via the official API, \S\ref{sec:closed-frontier-attempt}) finds defended ASR of 0.36, but the baseline-vs-defended difference is \emph{not statistically significant} ($p = 0.393$, $N=25$): this is consistent with the degradation hypothesis but under-powered and not confirmatory. The V2 structural defense's efficacy is LLM-independent by construction.
\end{itemize}

\paragraph{Data Availability}
To facilitate future research and reproducibility, we will open-source our attack payloads, the evaluation harness, and the anonymized real-LLM evaluation logs upon publication.

The remainder of this paper is organized as follows. Section~\ref{sec:threat-model} presents the threat model and formalizes the recovery-channel injection definition. Section~\ref{sec:defense} details the defense design and implementation. Section~\ref{sec:evaluation} reports the evaluation. Section~\ref{sec:conclusion} concludes.
